\section{A measured degree estimate}\label{sec:degree}

This section bounds the degree of a map from a measured family of
Riemannian manifolds to a cube.  Its integrand combines the weighted scalar
quantity with a separate allowance for each coordinate gradient.  The conclusion involves
ordinary Riemannian volume, integrated also in the transverse direction.
The family need not consist of complete manifolds, and its degree need not
be an integer.  We first specify the control at the sides of the cube that
makes this statement, and the integrations in its proof, well defined.

\subsection{Measured bands and degree}

We use the following model of a measured space of leaf manifolds.  The
underlying space $X$ is standard Borel and has a countable collection of
flow boxes
\[
  \kappa_a:T_a\times Q_a\longrightarrow X,
\]
whose images cover $X$ and which are Borel isomorphisms onto their images.
Here $Q_a$ is an open box
in $\mathbb R^d$, $T_a$ is standard Borel, and $T_a$ carries a
$\sigma$-finite measure $\mu_a$.  The plaques $\kappa_a(\{t\}\times Q_a)$
fit together as oriented smooth $d$-manifolds, each Hausdorff and second
countable, and each plaque is open in its leaf.  The overlap of two boxes is a countable union of pieces on
which the change of transverse coordinate is a Borel bijection between
Borel subsets, preserves the transverse measures, and the change of leaf
coordinate is a smooth orientation-preserving diffeomorphism on open
sets, depending Borel measurably on the transverse coordinate.  All
functions and tensor fields below are Borel in the boxes and smooth on
plaques.  This convention also applies to an open subset along the leaves,
whose leaves are the connected components of that subset.

The measures $d\mu_a(t)\,dV_h$ in the boxes define a single measured
Riemannian volume, denoted simply by $dV_h$.  The same convention, with
leafwise area in place of volume, defines $dA_h$ on regular submanifolds.
One may partition among countably many boxes to define either integral;
the equality on overlaps follows from invariance of the transverse
measures.  From a plaque, all plaques on its leaf have a countable Borel
enumeration, obtained by following finite strings of overlap maps.  In
particular the saturation of a transverse null set is still null in every
transversal.  We shall discard such saturations when an assertion holds on
almost all plaques.

For $t>0$ write $Q_t^d=(-t,t)^d$; in dimension zero this denotes one point.
We write $|G|_\infty=\max_i|G_i|$, with the empty maximum interpreted as
zero.

\begin{definition}[Controlled measured band]\label{def:degree-band}
Let $T>0$.  A controlled measured $d$-band over $Q_T^d$ consists of a
measured space of oriented $d$-manifolds as above, a leafwise metric $h$,
a function $\phi$, and a map $G=(G_1,\ldots,G_d):X\to Q_T^d$.
For every $0<t<T$ the part $\{|G|_\infty\le t\}$ has a countable flow-box
cover with the following properties.
\begin{enumerate}
\item The boxes have one common positive size and fixed positive buffers:
their smaller concentric boxes cover the indicated part.  They can be
chosen inside $\{|G|_\infty<t'\}$ for some $t<t'<T$.  After rescaling
coordinates we may, and do, call them unit boxes with buffers.
\item In these boxes the coefficients of $h$ and $h^{-1}$ and all their
leafwise coordinate derivatives are bounded by constants common to the
cover.  The same holds for all derivatives of $G$ and for all
positive-order derivatives of $\phi$.  No bound on the value of $\phi$ is
required.
\item The sum of the measures of the transversals of this cover is finite.
There is a common bound on the number of plaques of the cover whose
buffered boxes meet any one plaque of the cover.
\end{enumerate}
The cover and every bound in this definition may depend on $t$ and on the
given band.  The leaves are not assumed complete, and $G$ is not assumed
proper.
\end{definition}

For $d\ge1$, the stated control implies the version of Stokes' theorem used here.  To
see this, choose a smooth cutoff in each buffered box which is one on its
inner box.  Dividing the cutoffs by their sum gives a leafwise smooth,
Borel partition of unity on a slightly smaller inner band.  The overlap
bound makes the sum locally finite and bounds all its derivatives.  The
metric and inverse-metric bounds give uniform bounds for these derivatives
in geometric norms.  Consequently, for a leafwise smooth $(d-1)$-form
$\eta$ supported over a compact subcube, with $\eta$ and $d\eta$ uniformly
bounded there, chartwise Stokes and Fubini give
\begin{equation}\label{eq:degree-stokes}
                  \int_X d\eta=0.
\end{equation}
Indeed, each cutoff multiple of $\eta$ is compactly supported in its
plaque, and the sum of the absolute integrals is bounded by a constant
times the sum of the transverse measures.  The terms differentiating the
partition cancel because its sum is one.  The same argument gives measured
Stokes for leafwise currents with uniform local mass and boundary-mass
bounds in those boxes: apply the current identity in each plaque to the
cutoff multiple of a test form and then sum.  This includes test forms
whose pullbacks have noncompact support in a leaf.  We will always verify
the indicated local bounds before using current Stokes.

If $\omega$ is a smooth, compactly supported $d$-form on $Q_T^d$ with
$\int_{Q_T^d}\omega=1$, define
\begin{equation}\label{eq:degree-definition}
              \mathcal D_X(G)=\int_X G^*\omega.
\end{equation}
The integral is absolute: its support lies in an inner band, the
derivatives of $G$ there are bounded, and that band has finite measured
volume.  For $d\ge1$, a compactly supported top form on a cube of integral
zero has a compactly supported primitive.  Applying \eqref{eq:degree-stokes} to the
pullback of this primitive shows that \eqref{eq:degree-definition} is
independent of $\omega$.  This is the \emph{measured degree}, with the
usual signed orientation convention.  For $d=0$ it means the measured
signed count of the oriented points.  We abbreviate it to $\mathcal D$
when the band and map are understood.

\begin{theorem}[Measured degree estimate]\label{thm:degree}
There are absolute constants $C_0,C_1>0$ with the following property.
Let $q\ge2$, let $d$ be an integer with $0\le d\le q$, and let $\tau>0$.
For $0\le r\le d$ set
\begin{equation}\label{eq:degree-widths}
 T_r=2\tau+\frac{r\tau}{q},\qquad
 b_r=\tau\left(1+\frac1{2q}+\frac r{4q^2}\right).
\end{equation}
Suppose $(X,h,\phi,G)$ is a controlled measured $d$-band over $Q_{T_d}^d$
in the sense of Definition~\ref{def:degree-band}.  Let $B\ge1$ and
$L_1,\ldots,L_d>0$.  Suppose throughout $X$ that
\begin{equation}\label{eq:degree-horizon-bound}
             L_i|dG_i|_h\le B-\frac dq\qquad(1\le i\le d),
\end{equation}
and suppose that
\begin{equation}\label{eq:degree-horizon}
       \frac{C_0(1+B)^{2/3}}{L_i}<\frac{\tau}{10q}
                       \qquad(1\le i\le d).
\end{equation}
Define $H_i(G)=2(L_i|dG_i|_h-1)_+$.  Then
\begin{equation}\label{eq:degree-estimate}
 |\mathcal D_X(G)|(2\tau)^d
 \le \frac{\exp(C_1d/q)}{\prod_{i=1}^dL_i}
     \int_{\{|G|_\infty<b_d\}}
       \exp\left(\frac{-\wscal(h,\phi)+\sum_{i=1}^dH_i(G)}2\right)dV_h.
\end{equation}
For $d=0$ the product is one, the sum is zero, and
$\wscal(h,\phi)=0$; the assertion is the bound of a signed count by the
unsigned count.
\end{theorem}

The proof removes one coordinate at a time.  For $r\ge1$, the gaps
$T_r-T_{r-1}=\tau/q$ and $b_r-b_{r-1}=\tau/(4q^2)$ supply, respectively,
chart buffers and room for the residual value error when locating the
integration core, while the permitted scaled gradient of each retained
coordinate increases by $1/q$.  A family of boundaries
ordered by a real parameter supplies the factor $2\tau$ on the left.
Its normal speed supplies the coarea factor on the right.  In dimensions
where the boundaries can be singular we will construct a regular band
of the same degree; a scalar gain will pay exactly for the volume of the
graphs used to construct it.  The following formulation separates that
construction from the final induction.

\begin{lemma}[One step of measured descent]\label{lem:degree-descent}
Assume the hypotheses of Theorem~\ref{thm:degree} with $d\ge1$, put
$p=d-1$, and put
\[
       \Theta=\exp\left(\frac{-\wscal(h,\phi)+\sum_{i=1}^dH_i(G)}2\right).
\]
There is a Borel family of oriented regular hypersurface portions
$\Sigma_s$ in the leaves, $-\tau<s<\tau$, and, for almost every $s$, a
positive function $v_s$ on $\Sigma_s$, smooth along it and Borel in its
charts, with the following properties.  Transverse null subsets and their
saturations may be removed for each such $s$.
\begin{enumerate}
\item For every nonnegative Borel function $F$ supported in
$\{|G|_\infty<b_d\}$,
\begin{equation}\label{eq:degree-coarea}
 \int_{-\tau}^{\tau}\!ds\int_{\Sigma_s}v_sF\,dA_h
                       \le\int_XF\,dV_h.
\end{equation}
\item For each such $s$ there is a controlled measured $p$-band
$(X_s^+,h_s^+,\phi_s^+,J_s^+)$ over $Q_{T_p}^p$ of degree
$\mathcal D_X(G)$.  Its coordinates are indexed by $j=2,\ldots,d$ and
satisfy
\[
             L_j|dJ_{s,j}^+|_{h_s^+}\le B-\frac pq.
\]
With $H_j^+=2(L_j|dJ_{s,j}^+|_{h_s^+}-1)_+$, there is an absolute constant
$C_*>0$ such that
\begin{equation}\label{eq:degree-child-integral}
 \begin{split}
 &\int_{\{|J_s^+|_\infty<b_p\}}
        \exp\left(\frac{-\wscal(h_s^+,\phi_s^+)+\sum_{j=2}^dH_j^+}{2}\right)
        dV_{h_s^+}\\
 &\hspace{25mm}\le \frac{e^{C_*/q}}{L_1}
       \int_{\Sigma_s\cap\{|G|_\infty<b_d\}}v_s\Theta\,dA_h.
 \end{split}
\end{equation}
In dimension $p=0$ the left side is the measured unsigned count.
\end{enumerate}
The child in the second assertion may be chosen separately for each
typical $s$; no measurability of those choices in $s$ is asserted or needed.
\end{lemma}

\subsection{An ordered family of boundaries}

We prove Lemma~\ref{lem:degree-descent} in this and the next two
subsections.  Fix intermediate widths
\begin{equation}\label{eq:degree-intermediate-widths}
              T_p<t_1<t_2<t_3<t_4<T_d.
\end{equation}
Whenever a chart buffer is needed, choose further intermediate widths
between these.  The width $t_4$ controls confinement and the parent
charts, $t_3$ contains the boundary portions to be used, $t_2$ will bound
an auxiliary complete region, and $t_1$ will bound the child.  Such
choices can depend on the given band.

Choose once and for all a smooth even function $\gamma$ with
$0\le\gamma(y)\le y^2$ and with $\gamma(y)$ bounded above and below by
positive absolute multiples of $|y|^{3/2}$ for $|y|\ge1$.  Let $P_0$ be the
solution about zero of
\begin{equation}\label{eq:degree-pressure}
 P_0(0)=0,\qquad P_0'=L_1\psi(P_0),\qquad
       \psi(y)=1+\frac{\gamma(y)}{8(1+B)}.
\end{equation}
It is odd and increasing and has poles at $\pm b$, where
\[
 b=\frac1{L_1}\int_0^\infty\frac{dy}{\psi(y)}
       \le\frac{C(1+B)^{2/3}}{L_1}.
\]
For the inequality, split the integral at $(1+B)^{2/3}$ and use the
$|y|^{3/2}$ bound on the tail.  Choose the absolute $C_0$ in
\eqref{eq:degree-horizon} at least this $C$.  Thus
\begin{equation}\label{eq:degree-pole-width}
                         b<\frac{\tau}{10q}.
\end{equation}

\paragraph{The local slice calculation.}
We first record the analytic purpose of the pressure law.  Let $\Sigma$
be a smooth cooriented local hypersurface, with unit normal $\nu$, and set
\[
 \begin{gathered}
 D=H+\partial_\nu\phi,\qquad
 H=\operatorname{tr}\mathrm{II},\qquad
 \mathrm{II}(Y,Z)=\langle\nabla_Y\nu,Z\rangle,\\
 \mathcal P=|\mathrm{II}|^2+
             (\operatorname{Ric}_h-\nabla^2\phi)(\nu,\nu).
 \end{gathered}
\]
Put $\mathfrak a=L_1\psi(D)>0$.  When $D=P_0(s-G_1)$, the derivatives
of the raw pressure are $\partial_sP_0(s-G_1)=\mathfrak a$ and
$d_xP_0(s-G_1)=-\mathfrak a\,dG_1$.  The following calculation uses
only a positive smooth function $v$ satisfying the local equation
\begin{equation}\label{eq:degree-local-lapse}
 -\bigl(\Delta_{\Sigma,\phi}+\mathcal P
                 -\mathfrak a\,dG_1(\nu)\bigr)v=\mathfrak a.
\end{equation}
Here $\Delta_{\Sigma,\phi}$ is weighted by
$e^{\phi|_\Sigma}dA_h$, and slice norms use the induced metric.  The
ordered construction below will produce this equation on almost every
regular slice.

Put $\phi_\Sigma=\phi|_\Sigma+\ln v$ and
$S_\Sigma=\wscal(h|_{T\Sigma},\phi_\Sigma)$.  Gauss' equation and the
normal/tangential decomposition of derivatives of $\phi$ give
\[
 \wscal(h|_{T\Sigma},\phi|_\Sigma)
       =\wscal(h,\phi)+D^2-2\mathcal P+|\mathrm{II}|^2.
\]
Adding $\ln v$ to the weight changes this by
$-2\Delta_{\Sigma,\phi}v/v+|d\ln v|^2$.  Substitution of
\eqref{eq:degree-local-lapse} yields the exact identity
\begin{equation}\label{eq:degree-slice-identity}
 S_\Sigma=\wscal(h,\phi)+D^2-2\mathfrak a\,dG_1(\nu)
                 +|\mathrm{II}|^2+\frac{2\mathfrak a}{v}+|d\ln v|^2.
\end{equation}
This is also the hypersurface calculation behind
Lemma~\ref{lem:graph-identity}; it is valid for $p=0$ with left side zero.

Set $t=L_1|dG_1|_h\le B$ and $x=v/L_1$.  Since $\gamma(D)\le D^2$,
\[
 -2\mathfrak a\,dG_1(\nu)
       \ge-2t\left(1+\frac{\gamma(D)}{8(1+B)}\right),
 \qquad \frac{t\gamma(D)}{4(1+B)}\le\frac{D^2}{4}.
\]
We may discard the remaining nonnegative part of $D^2$, and use
$2\mathfrak a/v\ge2/x$.  Finally,
$t\le1+(t-1)_+$ and $x^{-1}+\ln x\ge1$ for $x>0$.  Thus
\eqref{eq:degree-slice-identity} implies
\begin{equation}\label{eq:degree-slice-bound}
 S_\Sigma\ge\wscal(h,\phi)-H_1(G)-2\ln(v/L_1)+|d\ln v|^2.
\end{equation}
The square term will be used later when the weight is changed near
singular points.  For $j=2,\ldots,d$, put
$H_j^\Sigma=2(L_j|d_\Sigma G_j|-1)_+$.  Tangential restriction gives
$H_j^\Sigma\le H_j(G)$, so the full pointwise density comparison is
\begin{equation}\label{eq:degree-raw-integrand}
 \exp\left(\frac{-S_\Sigma+\sum_{j=2}^dH_j^\Sigma}{2}\right)
 \le\frac v{L_1}e^{-|d\ln v|^2/2}\Theta
 \le\frac v{L_1}\Theta.
\end{equation}
For $p=0$ the sum and slice gradient are zero, and the left side is one.
The remaining work supplies the slices and speeds, and verifies the
controlled-band and degree conditions needed to use this comparison.

For a small $\eta>0$, prolong $P_0$ smoothly and strictly increasingly
past the points $\pm(b-\eta)$, keeping it unchanged on
$[-b+\eta,b-\eta]$.  Denote the resulting smooth function by
$\widehat P$ and set
\[
                \mathfrak h_s(x)=\widehat P(s-G_1(x)),\qquad -\tau<s<\tau.
\]
We will choose $\eta$ sufficiently small for the given band and then keep
it fixed.  On every boundary portion used below, $\widehat P$ will agree
with $P_0$ and hence obey \eqref{eq:degree-pressure}.

\paragraph{Least minimizers, including on noncompact leaves.}
Write $P_\phi(E;U)=\int_{\partial^*E\cap U}e^\phi\,dA_h$ for the weighted
perimeter of a set of locally finite perimeter.  We use local differences
of the functional
\begin{equation}\label{eq:degree-functional}
                  P_\phi(E)-\int_E\mathfrak h_s e^\phi\,dV_h.
\end{equation}
This is a weighted $\mu$-bubble functional.  For soap bubbles in Gromov's
terminology, see \cite[Section~5$\frac56$]{Gromov1996}; for the weighted
functional and its local first- and second-variation formulas, see also
\cite[Section~3, Equation~(1), and Lemmas~13--14]{ChodoshLi2024}.
Neither term is required to be finite on the whole leaf.  In particular,
this construction uses no global weighted-volume bound.

The output will be a canonical nested family $E_s$ on each open leaf,
independent of the exhaustion and jointly Borel in the plaque data and
$s$.  Each $E_s$ will minimize these local differences for compact
variations.  The local comparison argument also supplies the perimeter
compactness needed later to turn convergence of sets into graphical
convergence at regular limit points.  We establish these properties in
turn, before confining the boundaries to the raw-pressure region.

\emph{Compact domains and leastness.}
For a compact subset $K$ of a leaf, minimize \eqref{eq:degree-functional}
over sets which vanish off $K$, using their full perimeter in the leaf.
The empty set is a competitor.  On a neighborhood of $K$ the density is
smooth and positive and the forcing is bounded.  A bound on the functional
therefore gives a perimeter bound, and BV compactness and lower
semicontinuity give a minimizer; see, for example,
\cite[Theorem~12.26]{Maggi2012} for the compactness theorem.  Among the
minimizers choose one of minimum
ordinary volume.  The weighted perimeter satisfies
\[
 P_\phi(E\cap F)+P_\phi(E\cup F)\le P_\phi(E)+P_\phi(F),
\]
and the volume term for a fixed forcing is additive under union and
intersection.  Thus the intersection of any two minimizers is again a
minimizer.  Minimum volume forces the selected minimizer to be contained,
up to null sets, in every other one.  It is consequently unique as an
$L^1$ set; call it the least minimizer in $K$.

\emph{Order and exhaustion.}
The same inequality proves both monotonicities we need.  If $K\subset K'$
and two forcings satisfy $\mathfrak h\le\mathfrak h'$, a minimizer for
$(K,\mathfrak h)$ and one for $(K',\mathfrak h')$ obey
\[
 \mathcal F_{\mathfrak h}(E\cap F)+\mathcal F_{\mathfrak h'}(E\cup F)
              \le \mathcal F_{\mathfrak h}(E)+\mathcal F_{\mathfrak h'}(F),
\]
where $\mathcal F$ denotes \eqref{eq:degree-functional}; the inequality for
the volume terms follows on $E\setminus F$ from
$\mathfrak h\le\mathfrak h'$.  Minimality in
the two domains makes both inequalities equalities, and leastness in $K$
implies $E\subset F$.  Exhaust a leaf by increasing finite unions of
compact coordinate subboxes whose interiors exhaust it.  The least
minimizers increase with this exhaustion.  Their union, interpreted in
$L^1_{\mathrm{loc}}$, is independent of the exhaustion, since each compact
subset of either exhaustion is eventually contained in one of the other.
This defines $E_s$, and the forcing monotonicity makes $E_s\subset E_t$
for $s<t$.

\emph{Local comparison.}
The limit must still minimize for compact variations, since all later
regularity and compactness arguments use that property.
Ball replacement in a relatively compact chart gives uniform local
perimeter bounds for the exhaustion minimizers: choose a radius with
controlled slicing area, replace the set inside that ball by the full or
empty ball, and bound the forcing term by its local supremum times the
ball volume.  Thus the increasing sets converge locally in BV weakly and
in $L^1$.  To pass a comparison to the limit, surround the compact region
to be changed by a thin outer belt on which the proposed competitor
agrees with $E_s$.  Choose a Lipschitz coordinate level function across
the belt.  Coarea bounds the integral, over its levels, of the trace
mismatch between an exhaustion minimizer and $E_s$ by their $L^1$
mismatch in the belt.  One may therefore splice the proposed competitor
inside a level to the exhaustion minimizer outside it so that the extra
trace area tends to zero.  The levels can also be chosen with zero
limiting perimeter on their boundaries, since this excludes only a null
set of levels.  Minimality before passage to the limit, convergence of
the bounded forcing integrals, and lower semicontinuity give the compact
comparison for $E_s$.  The same argument applies under smooth local
convergence of the metric, density, and forcing coefficients.  Below we
apply it with the unchanged limit as competitor to prove convergence of
the local perimeter measures.  Multiplying a density by a positive
constant has no effect on any of these comparisons.

\emph{Borel dependence.}
We explain why all these choices are Borel, including jointly in $s$.
For this coding, subdivide the original atlas into bounded coordinate
boxes with compact closure in the older boxes, and affinely normalize
each to one fixed cube $Q=(-1,1)^d$.  Fix once a countable list
$(R_\ell)_{\ell\ge1}$ of all open rational boxes with closure in $Q$.
Starting at a plaque, enumerate its accessible normalized plaque
occurrences Borel measurably in slots $j\ge1$, using dummy slots when
necessary.  Use the fixed Polish space
\[
             \mathcal Z=\prod_{j,\ell\ge1}L^1(R_\ell),
\]
with Lebesgue $L^1$ norms and the product topology.  A set is coded by its
indicator on $R_\ell$ in occurrence $j$; dummy entries are zero.
Compatibility is tested on compact exhaustions of all overlap domains,
including repeated occurrences.  The closures of the first finitely
many accessible $R_\ell$ in a fixed enumeration of the pairs $(j,\ell)$
give the compact exhaustion domains.  The conditions that the code is an
indicator, agrees on overlaps, and vanishes off the specified domain are
Borel, and for fixed plaque data are closed in $\mathcal Z$.

Integrals against chartwise Borel smooth densities are Borel by Fubini.
Perimeter is Borel as well: its supremum over divergence integrals may
be taken over fixed countable families of rational smooth test fields in
the $R_\ell$, and finite sums of them.  Partitions of unity in a leaf
show density among all compactly supported tests.  Norm constraints are
checked on countably many dense coordinate points; overlap maps and a
partition by the first covering chart make the integrals Borel.

Fix one plaque-data fiber and one compact exhaustion domain $K$.  A
functional sublevel for sets supported in $K$ bounds their weighted
perimeter, since the forcing integral is uniformly bounded on this
fixed compact set.  In every active $R_\ell$ the metric and density
have a positive lower bound on its closure, so this gives a Euclidean BV
bound there; the indicator's $L^1(R_\ell)$ norm is bounded by $|R_\ell|$.
Diagonal BV compactness gives convergence in $\mathcal Z$ along a
subsequence.  The fiberwise closed indicator, overlap, and support
conditions persist.  The volume term is continuous, since $K$ is
covered by finitely many compact subboxes, and perimeter is lower
semicontinuous.  Each sublevel is therefore compact in this fixed Polish
product.  These compactness bounds may depend on the fiber and on
$R_\ell$; the projection theorem requires compact fibers, not a common
bound across plaque data.  Strict sublevels are countable unions of such
compact sections.  The Borel projection theorem for sets with
$\sigma$-compact sections therefore makes the minimum value a Borel
function of the chart data.  Apply it again to minimize volume in the
compact set of minimizers.  The selected section is a singleton in the
$L^1$ coding, by leastness, so the Borel singleton-section theorem gives
the selected minimizer as a Borel function.  These are standard forms of
the Arsenin--Kunugui and Lusin--Souslin theorems; see
\cite[Theorem~18.18 and Corollary~15.2]{Kechris1995}.  The exhaustion limits agree
for different starting plaques by exhaustion independence.  They are therefore Borel
on the measured space.  We may use the density-one representative of
$E_s$; the support of its boundary measure is also Borel, by testing its
mass on a countable basis of coordinate balls.

\paragraph{Confinement.}
For sufficiently small $\eta$, uniformly for $-\tau<s<\tau$, the boundary
support in $\{|G|_\infty\le t_4\}$ satisfies
\begin{equation}\label{eq:degree-confinement}
\operatorname{supp}\partial[E_s]\cap\{|G|_\infty\le t_4\}
                         \subset\{|s-G_1|<b-\eta\}.
\end{equation}
In this inner band the side $G_1\le s-b+\eta$ is full and the side
$G_1\ge s+b-\eta$ is empty, up to null sets.  To prove this uniform claim,
the ODE at the positive pole gives
\[
 b-x=\frac1{L_1}\int_{P_0(x)}^\infty\frac{dy}{\psi(y)}
                  \asymp_{\mathrm{inst}}P_0(x)^{-1/2}
                  \quad(x\uparrow b).
\]
Thus at distance comparable to $\eta$ from a pole the absolute forcing is
at least $c_{\mathrm{inst}}\eta^{-2}$.  The prolongation is monotone, so
this remains a lower bound beyond that position.  By the uniform
derivative bound on $G_1$, a chart ball of radius
$r=c'_{\mathrm{inst}}\eta$ centered at a point with $s-G_1\ge b-\eta$
has forcing at least $H=c''_{\mathrm{inst}}\eta^{-2}$ throughout it, if
$c'_{\mathrm{inst}}$ is small.  The balls are taken in a buffered atlas
for a slightly larger band than $t_4$.

Here is the phase-exclusion estimate in such a ball.  Normalize the
weighted density by its value at the center; the derivative bounds for
$\phi$ make this density uniformly comparable to coordinate volume in
the ball.  Let $V(t)$ be the normalized volume of the complementary
phase in a concentric coordinate ball $B_t$.  Comparison with the set
which fills $B_t$, and perimeter slicing, give for almost every $t<r$
\[
 P_{\mathrm{norm}}(E_s;B_t)+H V(t)\le C V'(t).
\]
The chart isoperimetric inequality applied to that complementary phase,
extended by zero outside $B_t$, gives, when $V(t)>0$,
\[
 V(t)^{(d-1)/d}\le C\bigl(P_{\mathrm{norm}}(E_s;B_t)+V'(t)\bigr).
\]
If the complementary phase occurs in $B_{r/4}$, the two inequalities first
give $V'\ge cV^{(d-1)/d}$ and hence $V(r/2)\ge cr^d$.  They also give
$V'\ge cHV$ on $[r/2,r]$.  It follows that
$V(r)\ge cr^d\exp(cHr)$, contrary to $V(r)\le Cr^d$ when $Hr$ is large.
This argument includes $d=1$, where the first differential inequality
has exponent zero.  Since $Hr\asymp_{\mathrm{inst}}\eta^{-1}$, small
$\eta$ excludes the complementary phase in the smaller ball.  Negative
forcing gives the empty phase by the same argument.  Applying these
balls beginning at $|s-G_1|=b-\eta$, with a small fixed extra margin in
the forcing estimate, proves \eqref{eq:degree-confinement}.  In particular
the unprolonged Equation~\eqref{eq:degree-pressure} holds on all the
boundary portions used in the proof.

\paragraph{The local regularity and compactness inputs.}
We record precisely the local geometric measure theory used from now on.
Set $\Sigma_s^{\mathrm{all}}=\operatorname{supp}\partial[E_s]$.  For
$d\ge2$, its regular portion is a smooth multiplicity-one hypersurface.
The remaining set is closed in $\Sigma_s^{\mathrm{all}}$ and has Hausdorff
dimension at most
\begin{equation}\label{eq:degree-singular-dimension}
                              p-7.
\end{equation}
It is empty if this number is negative.  These conclusions hold in each
relatively compact chart.  To place our functional under the standard
codimension-one theorem, normalize $\phi$ by a constant value $\phi_0$
in the chart and use the smooth metric
\[
                  \widehat h=e^{2(\phi-\phi_0)/p}h.
\]
Its area is the normalized weighted area.  Its volume forcing is
$\mathfrak h_s e^{-(\phi-\phi_0)/p}$, bounded smoothly in a smaller chart.
Minimality implies that its perimeter is minimizing up to a fixed
multiple of the volume of the symmetric difference for compact
variations.  The regularity theorem for such almost-minimizing boundaries
gives $C^{1,\alpha}$ regularity off a closed set of dimension
\eqref{eq:degree-singular-dimension}; see
\cite[Theorem A.1 and Remark A.3]{Eichmair2009}.  The weak prescribed
mean-curvature equation on this regular portion and elliptic regularity
then make it smooth.  The cited almost-minimizing formulation uses
integral-current variations.  It applies to these set variations as
follows.  Add a compactly supported ambient-dimensional integral current
to $[E_s]$ and truncate its integer BV multiplicity to the interval
$[0,1]$.  BV truncation does not increase boundary mass, and the volume
where the resulting set changes is at most the mass of the added current.
The set comparison therefore implies that formulation as well.

\emph{Local compactness claim.}
Suppose $d\ge2$ and $E_j$ minimize the absorbed-perimeter functionals
for all compact variations in a common buffered normalized outer chart.
Assume that their absorbed metrics and converted volume forcings
converge smoothly on smaller boxes and that $E_j\to E$ in
$L^1_{\mathrm{loc}}$.  Then $E$ is locally minimizing, the normalized
perimeter measures and the associated boundary varifolds converge
weakly locally to those of $E$, and the supports converge locally.
Near every regular point of the limiting boundary, the boundaries are
single smooth graphs converging in $C^\infty$ on a cylinder which may
depend on that point.  In particular, for every relatively compact
continuity subdomain $U$, meaning that the limiting perimeter of
$\partial U$ is zero, the perimeter masses on $U$ converge.  For the
minimizing boundaries considered here, every sequence of controlled
plaques from the fixed band and fixed forcing construction has a
subsequence satisfying these coefficient and set convergence hypotheses
after normalization.

Here is the proof, including the local perimeter assertion.  In the
common normalized coordinates subtract the value of $\phi$ at the chart
center.  The common bounds on all derivatives, and the metric and
inverse-metric bounds, make the normalized metric, weight, and forcing
coefficients precompact in $C^\infty$ on every smaller fixed box, by
diagonal Arzelà--Ascoli.  Pass to such a convergent subsequence; in
particular the absorbed metrics $\widehat h_j$ converge smoothly to a
positive metric $\widehat h_\infty$.  Ball comparison gives the local
perimeter bounds used above, so BV compactness gives a further
subsequence converging in $L^1_{\mathrm{loc}}$ to a set $E$.  The belt
comparison with any compact variation, as proved above and now with the
smoothly converging coefficients, shows that $E$ is locally minimizing.

Write $\mu_j=P_{\widehat h_j}(E_j;\,\cdot\,)$ and
$\mu_\infty=P_{\widehat h_\infty}(E;\,\cdot\,)$.  To prove convergence
of these measures, begin with a relatively compact coordinate domain
$V$ from a convergence-determining family, chosen so that
$\mu_\infty(\partial V)=0$.  Surround it by a thin outer belt with level
domains $V_t$ contained in an outer enlargement $V_\epsilon$.  Coarea
chooses a level $t_j$ at which the trace mismatch between $E_j$ and $E$
tends to zero.  Splice the unchanged $E$ inside $V_{t_j}$ to $E_j$
outside.  The common coefficient bounds make the extra perimeter at the
splice tend to zero, and the bounded converted forcing times the
$L^1$ mismatch tends to zero as well.  Minimality of $E_j$ gives
\[
       \limsup_{j\to\infty}\mu_j(V)\le\mu_\infty(V_\epsilon).
\]
Here the perimeter of the fixed set $E$ in the varying metric converges
by uniform convergence of the metric integrands on unit conormals.
Shrinking the belt makes the right side tend to $\mu_\infty(V)$, because
$\mu_\infty(\partial V)=0$.  Lower semicontinuity gives the reverse
inequality.  Such coordinate continuity domains determine weak local
convergence of Radon measures; the local mass bounds give compactness,
and the equality just proved identifies every subsequential limit.
Consequently $\mu_j$ converge weakly locally to $\mu_\infty$, and their
masses converge on every relatively compact continuity subdomain $U$.

For precision in passing to varifolds, restrict to a relatively compact
subdomain $U$ whose boundary has zero limit perimeter measure.  Such
continuity subdomains can exhaust the smaller box.  To check weak
convergence of the restrictions of $D\mathbf1_{E_j}$ as vector measures
on $\overline U$, extend a continuous vector test there to the outer
chart and multiply it by a cutoff supported in $U$ and equal to one away
from a $\delta$-collar of $\partial U$.  Local weak BV convergence
handles this cutoff test.  On the compact chart uniform ellipticity
gives $|D\mathbf1_{E_j}|\le C\mu_j$.  For a compact closed
$\delta$-collar $F_\delta$, weak perimeter convergence and the closed-set
inequality give
$\limsup_j\mu_j(F_\delta)\le\mu_\infty(F_\delta)$, which tends to zero
as $\delta\downarrow0$ because $\mu_\infty(\partial U)=0$.  The same
bound for the limiting vector measure lets the collar shrink and proves
the restriction convergence.  The perimeter masses on $U$ also converge.
In coordinates the perimeter
integrand for a metric $k$ is
$\mathfrak p_k(x,z)=\sqrt{\det k(x)}\,|z|_{k^{-1}(x)}$.
Uniform convergence of this integrand on unit conormals and the local
mass bound replace $\widehat h_j$ by $\widehat h_\infty$ in that mass
convergence.  This is convergence for one continuous positively
one-homogeneous integrand with Reshetnyak's strict convexity property:
equality in its triangle inequality occurs only for vectors on a common
nonnegative ray.  Reshetnyak's continuity theorem
\cite[Section~4, Theorem~3]{Reshetnyak1968} consequently gives convergence
for every continuous positively one-homogeneous integrand with linear
growth.  For a continuous compactly supported varifold test $a(x,P)$, use
$\mathfrak p_{\widehat h_\infty}(x,z)a(x,\ker z)$ for $z\ne0$ and zero
at $z=0$.  This gives precisely the local varifold convergence; the
uniform convergence of the metric integrands gives it also for the
original varying metrics.

To prove the support and graph assertions, the first variation in
the metric $\widehat h$ is, by BV Gauss--Green, integration of the bounded
volume forcing against the reduced-boundary conormal.  It therefore has
bounded mean curvature and no singular part.  The monotonicity formula
gives a common positive lower mass bound in a small ball about a support
point, starting with the unit density at almost every reduced-boundary
point.  This lower bound rules out support points a fixed distance from
the limit support.  Conversely, every ball about a point of the limit
support has positive limit mass, so weak convergence forces nearby
support points.  Near a fixed smooth multiplicity-one sheet of the
limit, Allard's interior regularity theorem gives single graphs with a
common $C^{1,\alpha_{\mathrm A}}$ bound for some
$\alpha_{\mathrm A}\in(0,1)$ on a smaller cylinder \cite{Allard1972}.
Compactness and identification of the already known varifold limit give
$C^{1,\beta}$ graph convergence for every $0<\beta<\alpha_{\mathrm A}$.
The prescribed equation, with the smoothly converging coefficients,
then gives $C^\infty$ convergence on still smaller cylinders by
interior elliptic estimates.  The cylinder here may depend on the fixed
regular point of the limit; uniform graph radii in later applications
will follow from their separate compactness contradictions.

The varying-metric use of these local theorems reduces to one smooth
ambient metric as follows.  Pass further in a smaller compact chart to a subsequence converging so
rapidly in $C^\infty$ that its metric is realized on slices $2^{-j}$ of
one extra coordinate, with the limiting metric on slice zero.  Indeed,
choose the error through order $j$ below $2^{-j^2}$ and interpolate by
disjoint smooth interval cutoffs.  The resulting product block-diagonal
metric is smooth; its slices have bounded second forms.  After smoothly
extending from a still smaller box to a closed manifold, a smooth Nash
embedding puts this situation in a single Euclidean ambient space
\cite[Theorem~2, p.~59]{Nash1956}.
The preceding monotonicity and Allard statements apply there, with the
same local bounded-mean-curvature property.  Applying this subsequence
argument to a sequence on which local convergence failed proves the
local support and graphical convergence asserted above.

For $d=1$ the needed facts are elementary.  A locally minimizing set has
locally finitely many jumps by the perimeter bound.  Two consecutive
jumps cannot be arbitrarily close in a controlled chart, since one can
replace the interval between them by the surrounding phase: delete it if
full, and fill it if empty.  This saves two weighted boundary costs
and changes the forcing integral by a quantity tending to zero.  Thus
the jumps have uniform local separation.  The same belt argument applies
directly to normalized weighted perimeter in dimension one, without
conformal absorption, and gives local perimeter convergence together
with BV convergence.  The separation then gives the corresponding
convergence of these zero-dimensional regular boundaries.  There are no
singular points in this case.

\subsection{The lapse, coarea, and the descended degree}

We now extract a smooth normal speed from the ordered family.  Write
$\Sigma_s$ for the regular part of $\Sigma_s^{\mathrm{all}}$ in
$\{|G|_\infty<t_3\}$.  The boundary has its exterior coorientation and
the induced orientation.  Its singular set has zero $p$-dimensional area,
so it never contributes to the integrals over the boundary current.

\paragraph{Differentiating the ordered family.}
On each leaf the sets $E_t$ converge to $E_s$ in $L^1_{\mathrm{loc}}$ as
$t\to s$ except at countably many $s$.  In fact, on a compact chart set
$K$ the function $t\mapsto\int_K\mathbf1_{E_t}$ is monotone, and nestedness
identifies its increments with the $L^1(K)$ distance between the
indicators.  It has at most countably many discontinuities.  Use the
countable chart exhaustion to obtain the assertion on the leaf.
At such a continuity parameter, the local convergence result above makes
$\Sigma_t$ converge smoothly and graphically near each regular point of
$\Sigma_s$ as $t\to s$.

In a smaller graph cylinder choose coordinates $(y,z)$ with the sets
below the graphs, and write the graph heights as $z=r(t,y)$.  They are
monotone in $t$.  They are strictly ordered wherever two such regular
graphs are defined: if $r(t,\cdot)\ge r(s,\cdot)$ for $t>s$ and they touch,
their tangent planes agree at contact.  The comparison of the graph mean
curvatures at that minimum of $r(t,\cdot)-r(s,\cdot)$ gives
$D_t\le D_s$, whereas their equations give
$D_t=\mathfrak h_t>\mathfrak h_s=D_s$ at the same point.  This is
impossible.  We use the conventions for $D$ and $\mathrm{II}$ from the
local slice calculation, with $\nu$ the exterior unit normal.  The
equation $D=\mathfrak h_s$ is the first variation of
\eqref{eq:degree-functional} on a regular portion.

Choose countably many possible graph cylinders in the chart enumerations,
and in each choose a countable dense set of $y$-coordinates.  At each of
these coordinates a monotone height has a finite derivative for almost
every parameter.  The test is meaningful even outside an interval of
graphical representation: use instead the monotone column function
\[
          t\longmapsto\int_{z_-}^{z_+}\mathbf1_{E_t}(y,z)\,dz,
\]
which equals $r(t,y)-z_-$ in the single-graph situation.  Consequently,
for almost every continuity parameter $s$, the nonnegative difference
quotients
\[
          w_t(y)=\frac{r(t,y)-r(s,y)}{t-s},\qquad t\ne s,
\]
converge to the finite column-derivative values as $t\to s$ at all these
dense test coordinates.  In particular they are bounded there.

In positive graph dimension the $w_t$ solve linear elliptic equations on
each smaller cylinder.  This follows by subtracting the two prescribed
graph equations and integrating their linearizations between the two
graphs.  Smooth graphical convergence gives common smooth coefficient
bounds and ellipticity; the explicit $t$-dependence of $\mathfrak h_t$ gives a
smooth bounded source.  Harnack's inequality with a bounded source,
applied using any one of the bounded test values, bounds $w_t$ on a
smaller patch.  Interior elliptic estimates then give smoothly convergent
subsequences there; these are the standard local Harnack and interior
estimates for uniformly elliptic equations with bounded lower-order
coefficients and source, as in \cite[Chapters~6 and~8]{GilbargTrudinger}.
Every such limit takes the already determined
derivative values at the dense test coordinates, so the limits agree.
Thus the whole difference quotient has a smooth local limit.  Its normal
component is a smooth function $v\ge0$ on the regular portion.  For
$p=0$ the same conclusion follows from the ordinary almost-everywhere
derivative of a monotone height.

Differentiation of $D=\mathfrak h_s$ now identifies this normal speed.
Use $\mathcal P$ from the local slice calculation and put
$\mathfrak a=\partial_s\mathfrak h_s=L_1\psi(D)>0$; the last equality
uses confinement, where the forcing agrees with the raw pressure.
The normal variation formula, with the tangential derivative of
$D-\mathfrak h_s=0$ cancelling, is
\begin{equation}\label{eq:degree-lapse}
 -\bigl(\Delta_{\Sigma,\phi}+\mathcal P+\partial_\nu\mathfrak h_s\bigr)v
                               =\mathfrak a.
\end{equation}
The convention for $\Delta_{\Sigma,\phi}$ is weighted by
$e^{\phi|_\Sigma}dA_h$.  The variation of $\partial_\nu\phi$ adds
$v\nabla^2\phi(\nu,\nu)-\langle d\phi,dv\rangle_\Sigma$ to the ordinary
mean-curvature variation, giving the signs in
\eqref{eq:degree-lapse}.  A nonnegative $v$ cannot vanish at an interior
point: there the left side is nonpositive, contrary to $\mathfrak a>0$.
Thus $v>0$.  In graph dimension zero the differential terms vanish and
the same equation rules out $v=0$.

These assertions hold jointly in the measured charts for almost every
$s$.  Here are the measurable details needed for that conclusion.
Single-graph regions with buffers can be recognized by countably many
cylinder tests on the $L^1$ subgraph codes, with slope and derivative
bounds and with room at the top and bottom.  Their smooth graph
representative is unique; the Borel singleton-section theorem applied to
its smooth-function code makes it Borel.  The continuity and
differentiability tests above use countable sequences, by monotonicity.
The resulting local derivatives and $v$ are Borel.  Fubini on each
$\sigma$-finite transversal now gives, for almost every $s$, a null set
of exceptional plaques.  There are countably many charts and overlap
moves, so their saturations can be discarded simultaneously.  In
particular the family of regular portions and the function $v$ used in
the integral below have the asserted Borel descriptions.
Set $v=0$ on the exceptional parameters and plaques when writing an
integral over the whole parameter interval.

In a graph cylinder, normal speed times area has the expression
\[
                v\,dA_h=J_h(y,r(s,y))\,\partial_s r(s,y)\,dy,
\]
where $J_h(y,z)\,dy\,dz$ is ambient Riemannian volume in the coordinates.
The one-dimensional change-of-variables inequality for a monotone
function says that integration against its almost-everywhere derivative
counts at most integration in $z$.  A jump or singular part of the
monotone function only removes mass from the former integral.  Apply this
to each column and then integrate in $y$ and the transversal.  Partition
the counted portions among the countably many graph cylinders so they
are not counted twice within one parameter.  They cannot give the same
ambient point at two parameters, by the noncontact argument above.
Therefore no ambient point is counted more than once on the right.
Equivalently, the counted images can be taken Borel by the Borel
one-to-one image theorem.  Transverse invariance on overlaps gives
\begin{equation}\label{eq:degree-coarea-proof}
          \int_{-\tau}^{\tau}\!ds\int_{\Sigma_s}vF\,dA_h
                               \le\int_XF\,dV_h
\end{equation}
for every nonnegative Borel $F$ supported in $\{|G|_\infty<t_3\}$.
Since $b_d<t_3$, this proves \eqref{eq:degree-coarea}.

We will require one more consequence of \eqref{eq:degree-lapse}.  On a
family of regular graph patches with common buffers and graph estimates,
all positive-order derivatives of $\ln v$ have uniform bounds on smaller
patches.  Indeed the coefficients in \eqref{eq:degree-lapse} and the
derivatives of $\mathfrak a$ are uniformly bounded there, and
$\mathfrak a$ has a uniform positive lower bound.  Confinement provides
the finite margin from the pressure poles needed for these bounds.  If
the value of $v$ could tend to zero at an interior point of such patches,
Harnack's inequality with bounded source would bound $v$ on a smaller
patch, and elliptic compactness would give a nonnegative limit with an
interior zero satisfying \eqref{eq:degree-lapse}, an impossibility.
Thus $v$ has a positive local lower bound.  Harnack then bounds its local
supremum by a constant times its value at an interior point, since the
bounded source can be absorbed using that lower bound.  Interior
estimates give $|\nabla^r v|\le C_r v$ on still smaller patches, and
differentiating $\ln v$ gives the claim.  This does not bound the value
of $\ln v$, which is not required by Definition~\ref{def:degree-band}.
For $p=0$ the assertion about positive-order derivatives is empty.

The local slice calculation now applies to the actual boundaries.
Confinement gives $D=\mathfrak h_s=P_0(s-G_1)$ and hence
$\partial_\nu\mathfrak h_s=-\mathfrak a\,dG_1(\nu)$ on every retained
regular portion.  Equation~\eqref{eq:degree-lapse} therefore becomes
exactly \eqref{eq:degree-local-lapse}.  With
$\phi_\Sigma=\phi|_\Sigma+\ln v$, the scalar bound
\eqref{eq:degree-slice-bound} and the pointwise comparison
\eqref{eq:degree-raw-integrand} hold on these portions.  The derivative
bounds just proved supply the required local controls for this weight.

\paragraph{Degree on the boundary.}
Let $J=(G_2,\ldots,G_d)$.  Fix a unit-integral one-form
$a_-(z)\,dz$ supported in $(-1.9\tau,-1.6\tau)$ and one
$a_+(z)\,dz$ supported in $(1.6\tau,1.9\tau)$.  Let $\beta$ be a
unit-integral smooth $p$-form supported in a sufficiently small cube
about zero in the residual target; for $p=0$ let $\beta=1$.  Set
\[
 \chi_{\mathrm{deg}}(z)=\int_{-\infty}^{z}(a_-(u)-a_+(u))\,du.
\]
This function is compactly supported in $(-1.9\tau,1.9\tau)$ and is one
between the two supports.  By \eqref{eq:degree-pole-width} and
\eqref{eq:degree-confinement}, the low support is in $E_s$, the high
support is outside $E_s$, and $\chi_{\mathrm{deg}}(G_1)=1$ on the confined
boundary where the residual form is supported.  Measured current Stokes therefore gives
\begin{equation}\label{eq:degree-boundary-degree}
 \begin{split}
 \int_{\partial[E_s]} \chi_{\mathrm{deg}}(G_1)J^*\beta
   &=\int_{E_s}(a_-(G_1)-a_+(G_1))\,dG_1\wedge J^*\beta\\
   &=\int_Xa_-(G_1)\,dG_1\wedge J^*\beta
     =\mathcal D_X(G).
 \end{split}
\end{equation}
All forms here are supported in an inner band.  The current $[E_s]$ has
bounded local volume and its boundary has the uniform local perimeter
bounds proved above, so this is precisely the current version of
\eqref{eq:degree-stokes}.  The singular set contributes no area, and
the counted part of the boundary is in $\{|G|_\infty<t_3\}$.
Consequently \eqref{eq:degree-boundary-degree} computes the degree by
$J^*\beta$ on that confined regular branch, with the induced orientation.
For $p=0$ it gives its measured signed count.

If $2\le d\le7$, this regular branch already provides the child.  To verify
the required uniformity, fix a typical $s$.  If graph radii and estimates
could not be chosen uniformly about the boundary in an inner region of
$\{|G|_\infty<t_3\}$, take offending points in buffered parent charts
and pass to normalized chart limits.  Local BV and varifold compactness
gives a limiting boundary through the limiting point.  There are no
singular points because $p-7<0$, so smooth graphical convergence near
that point supplies a common graph radius and all interior estimates,
contradicting the choice of points.  A fixed fine mesh of graph cylinders
now gives at most a bounded number of patches in each parent plaque.
Their transversals are the restricted parent transversals, counted with
the patch index.  Their total transverse measure and plaque overlap
remain bounded.  The graph estimates, together with the bounds for
derivatives of $\ln v$, give the band controls for the induced metric and
$\phi_\Sigma$.

Restrict this branch to
\[
                   |G|_\infty<t_1,\qquad |J|_\infty<T_p.
\]
The $G_1$ sides of the first restriction cannot be reached by confinement,
since $|G_1|<\tau+b<T_p$.  Near any other side some residual coordinate
has absolute value near $t_1>T_p$.  Hence every inner residual subcube
has buffered graph charts in this restriction.  Its degree equals
$\mathcal D_X(G)$ by \eqref{eq:degree-boundary-degree}, and its coordinate gradients only
decrease on restriction to the tangent space.  We take $J^+=J$ and use
\eqref{eq:degree-slice-bound}.

If $d=1$, so $p=0$, let $X_s^+$ be the confined regular points in
$\Sigma_s\cap\{|G_1|<t_1\}$, each with its boundary orientation, and let
$J^+$ be the unique map to the point $Q_{T_0}^0$.  Give these point leaves
the induced zero-dimensional metric and the weight $\phi_\Sigma$.
Confinement gives $|G_1|<\tau+b<T_0<t_1$ on the counted branch, so this
includes every point in \eqref{eq:degree-boundary-degree}.
The local jump separation proved above bounds the number of such points
in each buffered parent plaque.  Choose a fixed finite half-open mesh
finer than that separation in the normalized plaque.  Each mesh cell
contains at most one counted point.  The Borel boundary-support code and
the singleton-section theorem make that point a Borel function of the
transverse parameter wherever it exists.  These functions give a bounded
number of point charts, with transversals restricted from the parent and
counted with the mesh index.  Their total transverse measure is finite and their
plaque overlap remains bounded.  Point boxes have the required common
buffers, and all metric, map, and positive-order
weight derivative conditions in Definition~\ref{def:degree-band} are
vacuous in dimension zero, so this is a controlled $0$-band.  Its measured
signed count is $\mathcal D_X(G)$ by
\eqref{eq:degree-boundary-degree}; its integral is the measured unsigned
count.  At each point \eqref{eq:degree-raw-integrand} reads
$1\le(v/L_1)\Theta$, and confinement places every counted point in
$|G_1|<b_d$.  We return to the common integral estimate after treating
the possible singularities when $p\ge7$.

\subsection{A regular descendant in higher dimensions}

Assume now $p\ge7$, and fix one of the typical parameters $s$ above.
The singular set has zero $p$-dimensional area, so it does not change the
boundary-current integral.  Deleting it can, however, leave ends of the
regular locus over interior residual values, where regular chart radii
have no common lower bound.  We must keep those ends away from every
inner target subcube and preserve the degree computed by the closed
boundary current.

We will construct a controlled child carried by an open regular subset
of this slice, with the inherited orientation and transverse measure,
and with degree $\mathcal D_X(G)$.  For each residual row we require
$|dJ_j^+|_{h^+}\le |dG_j|_h+\varepsilon_{\mathrm{grad}}$ and the
one-sided value bound
$|J_j^+|\ge |G_j|-\varepsilon_{\mathrm{val}}$, with
\begin{equation}\label{eq:degree-tolerances}
 \sum_{j=2}^d L_j\varepsilon_{\mathrm{grad}}<\frac1{10q},
 \qquad \varepsilon_{\mathrm{val}}<b_d-b_p=\frac{\tau}{4q^2}.
\end{equation}
The gradient bound controls the residual penalties and the child's
coordinate bound.
The value bound, together with confinement of $G_1$, puts the child's
$b_p$ integration region over the parent's $b_d$ region.  Write
$\mathcal W=dV_{h^+}/dA_h$ for its ordinary leafwise volume density
relative to the original induced metric.  The scalar and volume target is
\begin{equation}\label{eq:degree-child-scalar}
 \begin{gathered}
 dV_{h^+}=\mathcal W\,dA_h,\\
 \wscal(h^+,\phi^+)\ge\wscal(h,\phi)-H_1(G)-2\ln(v/L_1)
                                      +2\ln\mathcal W-\frac1q.
 \end{gathered}
\end{equation}
The coefficient $2$ makes the volume factor cancel after exponentiation,
leaving the same lapse factor as in \eqref{eq:degree-raw-integrand}.
The density $\mathcal W$ includes neither $v$ nor the transverse measure.

Every constant and error chosen inside this subsection may depend on this
slice and on the band, but must be common across its retained plaques.
The choices proceed in three groups.  First a small parent perturbation
creates a target gap around the points without uniform regular charts.
Next we fix the kernel and graph constants using the controlled parent
data and that gap.  Finally we choose the marked-set cover small enough
for all scalar errors.  Its positive radius lower bound controls the
number and derivatives of the kernels for the graph solves; it does not
enter the scalar constants.  The adjustable scalar error in
\eqref{eq:degree-child-scalar} is at most $1/q$, so chart and elliptic
constants will not enter the absolute loss in Theorem~\ref{thm:degree}.

Choose a parent atlas with buffers around $\{|G|_\infty\le t_3\}$ still
inside $\{|G|_\infty<t_4\}$.  One obtains it from an atlas for a larger
inner cube by a bounded mesh subdivision, restricting transversals by
Borel tests as necessary.  Its plaque overlaps, also for the enlarged
supports below, remain bounded.  All operations in the parent made below
use this atlas.

\paragraph{Uniform regularity tests and avoidance of the remaining points.}
We use finitely many graph tests at a time to obtain a regular region
with common chart controls.  The points not yet recognized will be
excluded from the child, and the target gap below will keep the later
homotopies on the regular locus away from them.
Use one common normalized outer cube $Q$ for the buffered parent plaques.
Fix in $Q$ one ordered countable list of tests for graph cylinders, using
rational centers, radii, bounds, and a fixed countable dense set of
rotations.  Transport this same list, in this same order, by every parent
plaque chart; the list is independent of the plaque and its transverse
parameter.  A test recognizes that the boundary
in its outer cylinder is exactly one smooth graph, with a specified slope
bound and strict room at its top, bottom, and sides, and that the set is
on its specified side.  Its inner cylinder is then an open regular
region.  At level $k$ use the first $k$ tests of this fixed list in every
plaque, with their specified positive buffers.  Elliptic estimates on these
graphs give all interior derivative bounds for a fixed level.  Make the
tests with strict margins, so a regular limit point recognized by one of
them is recognized by that same test on every sufficiently close smoothly
convergent graphical boundary.  The graph coding already described shows
that recognition is Borel.

Let $D_k$ be the part of the boundary support in
$\{|G|_\infty\le t_3\}$ left after removing the open inner portions
recognized through level $k$ in any parent plaque.  These are decreasing
Borel sets, closed along each leaf.  Every regular point is eventually
removed.  More precisely, in any converging sequence of controlled
charts with levels tending to infinity, the $D_k$ can accumulate on a
compact subbox only at singular points of the limiting boundary.
Outside the limit support this follows from support convergence.  If
points of $D_{k_j}$ with $k_j\to\infty$ converged instead to a regular
limit point, one test with a fixed index $\ell$ in the common list would
recognize that point with strict room.  Smooth graphical convergence
makes the same test recognize the nearby points for large $j$, and then
$k_j\ge\ell$, contradicting their membership in $D_{k_j}$.  This is the
compactness fact used in both constructions that follow.

First make an arbitrarily small smooth perturbation $\bar J$ of $J$ in
the \emph{parent} band, small in value and first derivative, with all
higher derivatives uniformly bounded.  We claim it can be made so that,
for some $a_0>0$ and some finite $k$,
\begin{equation}\label{eq:degree-gap}
                |\bar J|_\infty>4a_0
             \quad\hbox{on a uniform neighborhood of }D_k.
\end{equation}
After the initial homotopy on the full boundary current, this gap will
keep the homotopies over a small target cube inside a uniformly regular
region.
To prove the local assertion behind the claim, fix a positive budget for
a constant shift in $\mathbb R^p$ and an inner compact plaque box.
There are a level and a positive gap, uniform across the controlled
plaques, such that some shift in that budget gives this gap from zero on
the marked set in the box.  Otherwise take a sequence of failing plaques,
with levels tending to infinity and requested gaps tending to zero.
Pass to normalized chart limits, including a smooth limit of the maps.
On a slightly larger compact box the marked sets concentrate at the
singular set of the limit.  By \eqref{eq:degree-singular-dimension} this
set has dimension at most $p-7$, so its image under the smooth limit map
has $p$-dimensional measure zero.  That image is compact.  A shift in the
given open budget can therefore avoid it with a positive distance.
Uniform convergence of the maps and concentration of the marked sets
contradict the alleged failure.  This proves the local uniform assertion.

Apply these shifts with smooth parent-chart cutoffs equal to one on the
inner boxes.  For clarity, simultaneous choices can be made Borel and
with bounded overlap.  The graph whose vertices are the buffered plaques
and whose edges record meeting supports is a Borel graph of finite
degree.  It has a countable Borel proper coloring: a countable separating
family of Borel sets distinguishes each vertex from its finitely many
neighbors, and the first finite distinguishing pattern is a countable
color.  The finite neighbor lists can be enumerated Borel measurably by
the Borel countable-section theorem.  Processing these color classes
successively and taking the first one of $\Delta+1$ colors avoided by
all earlier neighbors, where $\Delta$ bounds the degree,
gives a finite Borel proper coloring.  In each color the cutoff supports
are disjoint.  Choose rational shifts satisfying the local gap test; the
test on a compact inner box is Borel by projection with compact sections.
After a color is treated, increase the level if necessary and choose the
next shift budget smaller than the preceding positive gaps.  The marked
sets decrease, so earlier gaps persist.  There are only finitely many
colors.  The budgets can also make the total value and first-derivative
change as small as prescribed; the cutoffs give uniform bounds for all
higher derivatives.  Taking $a_0$ smaller than the final gaps and then
using the derivative bound of $\bar J$ extends the gap to a uniform
neighborhood.  This proves \eqref{eq:degree-gap}.

Second, set
\begin{equation}\label{eq:degree-alpha}
                         \alpha=p-\frac{13}{2}>0.
\end{equation}
The next cover will support a function large on the remaining marked
points while keeping the total kernel coefficients as small as needed.
For any prescribed positive budget and upper radius, a sufficiently
large $k'\ge k$ permits a cover of $D_{k'}$ in each inner plaque box by
finitely many small balls, of radii $r_i$, such that
\begin{equation}\label{eq:degree-cover-budget}
              \sum_i r_i^\alpha<\text{the prescribed budget}.
\end{equation}
The radii can also have a common positive lower bound, depending on the
budget and the controlled family.  To verify the uniform claim, suppose
it failed for levels tending to infinity and for requested lower radii
tending to zero.  Pass to a limiting chart.  Its compact singular set has
zero $\alpha$-dimensional Hausdorff measure because
$p-7<\alpha$.  Cover it on a slightly larger compact box by finitely
many balls with radii below the prescribed upper radius and with the
$\alpha$-sum strictly below the budget, allowing room to enlarge the balls.
This finite cover has a positive minimum radius.  The concentration
observation about $D_k$ makes the enlarged balls cover the marked points
of the sequence eventually, contradicting failure.  Rational chart balls
and fixed comparison factors give Riemannian balls with the same
conclusion.  Keep only balls meeting the marked set, allowing a small
change of centers and the reserved enlargement.  The centers then also
lie near $D_k$.  Such finite covers can be selected Borel: enumerate
finite tuples of rational balls, and test whether the compact marked
part outside their union is empty.  This is again a Borel compact-section
test.  The positive lower radius and \eqref{eq:degree-cover-budget}
bound the number of balls in any plaque.

The two uses of \eqref{eq:degree-alpha} are distinct: its strict excess
over $p-7$ made this cover possible, while
$\alpha+4=p-5/2<p-2$ will make the following potentials superharmonic.

\paragraph{Potentials near the marked set.}
Choose a small parent normal-coordinate radius $R>0$ so that a ball of
that radius about any cover center stays in the gap region
$|\bar J|_\infty>3a_0$ and in the buffered atlas.  It can also be chosen
small with respect to the local metric, weight, and forcing bounds.
This radius is fixed before the final cover is chosen; require its ball
radii to be much smaller than $R$.  For a center $c_i$ and radius $r_i$
let
\[
          K_i(x)=(r_i^2+\operatorname{dist}_h(x,c_i)^2)^{-1/2}
\]
on the smaller normal ball, and cut it off within radius $R$.  Choose a
radial cutoff which is one on the smaller ball, is positive before the
outer boundary, and near that boundary has the tail
$\exp(-1/(1-t))$ for $t<1$ and zero for $t\ge1$, where
$t=\operatorname{dist}_h(x,c_i)^2/R^2$.  Join this tail smoothly and
positively to the inner plateau.  Every positive power of the cutoff
then has a smooth exponential flat tail, so each power of $K_i$ used
below is smooth.  Denote its support by $V_i$.

On the regular boundary in $V_i$, for $j=\alpha,\alpha+4$,
\begin{equation}\label{eq:degree-kernel-bounds}
 \begin{split}
 -\Delta_{\Sigma,\phi}K_i^j&\ge cK_i^{j+2}-C\mathbf1_{V_i},\\
 |dK_i^j|&\le C\bigl(K_i^{j+1}+\mathbf1_{V_i}\bigr).
 \end{split}
\end{equation}
The constants may depend on the fixed local controls and on $p$, but are
independent of $r_i$ and of the eventual number of balls.  To check the
first estimate before cutoff, put $u=\operatorname{dist}_h(x,c_i)^2$.
The normal-coordinate Hessian bound, the bounded mean curvature
$H=\mathfrak h_s-\partial_\nu\phi$, and the bounded weight gradient give
\[
 \Delta_{\Sigma,\phi}u\ge2p-C\sqrt u,\qquad |du|_\Sigma^2\le4u.
\]
Differentiating $(r_i^2+u)^{-j/2}$ therefore gives
\[
 -\Delta_{\Sigma,\phi}K_i^j
 \ge j\left(p-j-2-C\sqrt u\right)K_i^{j+2}.
\]
Both chosen exponents are less than $p-2$, so small $R$ makes the
coefficient positive.  The gradient estimate follows by one
differentiation.  On the cutoff annulus the distance from the center is
bounded below by a fixed multiple of $R$, and all cutoff derivatives are
bounded independently of $r_i$; this accounts for the displayed constant
losses.  Outside $V_i$ they vanish, as the indicators record.

For $j=\alpha,\alpha+2,\alpha+4,\alpha+6$ put
\begin{equation}\label{eq:degree-potentials}
 u_j=\varepsilon_{\mathrm{base}}^j+\sum_i M_*r_i^\alpha K_i^j,
 \qquad \Lambda=\left(\frac{u_{\alpha+4}}{u_\alpha}\right)^{1/4},
 \qquad
 \mathcal L=|d\ln u_\alpha|^2+|d\ln u_{\alpha+4}|^2.
\end{equation}
These quantities have separate roles.  We will make $u_\alpha$ large on
$D_{k'}$ so that its logarithm locates the region to be excluded.
The term $\Lambda^2$ will supply the positive scalar contribution that
pays for the later coordinate graphs, while $\mathcal L$ controls the
logarithmic derivatives introduced by the weight and the first graph.
Here $0<\varepsilon_{\mathrm{base}}<1$ will be small; $M_*$ will be a
large finite number fixed before the final cover selection.  At a point, the sum is over the
kernels active there.  The overlap bound and
\eqref{eq:degree-cover-budget} allow the total active coefficient
$\sum_i M_*r_i^\alpha$ to be made arbitrarily small, uniformly at every
point.  In particular, for any $\delta>0$, first taking
$\varepsilon_{\mathrm{base}}$ small and then the cover budget small enough
gives, for $j=\alpha,\alpha+4$,
\begin{equation}\label{eq:degree-potential-laplacian}
       -\frac{\Delta_{\Sigma,\phi}u_j}{u_j}
                    \ge c\frac{u_{j+2}}{u_j}-\delta.
\end{equation}
Indeed \eqref{eq:degree-kernel-bounds} leaves the errors
$c\varepsilon_{\mathrm{base}}^{j+2}+C\sum_{i:x\in V_i}M_*r_i^\alpha$
before division by $u_j\ge\varepsilon_{\mathrm{base}}^j$.  Choose
$\varepsilon_{\mathrm{base}}^2$ and then the active sum divided by
$\varepsilon_{\mathrm{base}}^{\alpha+4}$ small enough for the displayed
$\delta$.

At each point the $u_j$ are the moments of positive atoms: one atom at
$\varepsilon_{\mathrm{base}}$ of mass one, and atoms at $K_i$ of masses
$M_*r_i^\alpha$.
Make their total active mass at most two.  Moment inequalities give
\begin{equation}\label{eq:degree-moments}
 \frac{u_{\alpha+6}}{u_{\alpha+4}}\ge\Lambda^2
     \ge\frac{u_{\alpha+2}}{u_\alpha},\qquad
 |d\ln\Lambda|^2\le\frac{\mathcal L}{8},\qquad
 |d\ln u_\alpha|\le C(\Lambda+1),\qquad
 \Lambda\ge c' u_\alpha^{1/\alpha}.
\end{equation}
For the first pair, use the probability weights proportional to the
atoms' masses times their $\alpha$-powers.  If $Z$ is the square of the
atom position, then
$\mathbb E Z\le(\mathbb E Z^2)^{1/2}$ and
$\mathbb E Z^3\ge(\mathbb E Z^2)^{3/2}$ give the two inequalities.  The
second estimate follows directly by differentiating
$\ln\Lambda=(\ln u_{\alpha+4}-\ln u_\alpha)/4$.  For the third, the
gradient bound in \eqref{eq:degree-kernel-bounds} bounds
$|du_\alpha|$ by $C\sum_iM_*r_i^\alpha(K_i^{\alpha+1}+1)$.
The first part divided by $u_\alpha$ is at most $\Lambda$ by H\"older's
inequality for those probability weights; the second is at most one after
the same small active-mass choice.  Finally, monotonicity of the
$L^r$ norms for a measure of mass at most two gives
$u_{\alpha+4}/u_\alpha\ge(u_\alpha/2)^{4/\alpha}$, proving the last
inequality.  These constants are independent of $M_*$ and $k'$.
At a point of $D_{k'}$ some ball has distance at most its radius, so its
term in $u_\alpha$ is at least $c''M_*$.  Hence
\begin{equation}\label{eq:degree-potential-on-bad}
                        u_\alpha\ge c''M_*\quad\hbox{on }D_{k'}.
\end{equation}
The positive lower bound on all radii and the bounded number of balls per
plaque make all $u_j$ smooth in the parent with uniform derivative bounds
for the final slice.  These bounds are allowed to depend on that radius
lower bound.

Add $\eta_0=(\ln u_\alpha+\ln u_{\alpha+4})/4$ to $\phi_\Sigma$ on the
regular boundary.  Put $Q_1=d\eta_0$.  The weight-change formula used in
\eqref{eq:degree-slice-identity}, now with
$\Delta_{\Sigma,\phi_\Sigma}=\Delta_{\Sigma,\phi}
+\langle d\ln v,d\,\cdot\,\rangle$, yields from
\eqref{eq:degree-potential-laplacian} the scalar gains
$c\Lambda^2+\mathcal L/2-2\langle d\ln v,Q_1\rangle-|Q_1|^2-\delta$.
Here the high moment in \eqref{eq:degree-moments} supplies $\Lambda^2$.
Since $|Q_1|^2\le\mathcal L/8$,
\[
 \frac{\mathcal L}{2}+|d\ln v|^2
       -2\langle d\ln v,Q_1\rangle-|Q_1|^2
       \ge\frac{\mathcal L}{2}-2|Q_1|^2\ge\frac{\mathcal L}{4}.
\]
Combining with \eqref{eq:degree-slice-bound}, and decreasing $c>0$ if
necessary, gives for the new weight
\begin{equation}\label{eq:degree-weight-gain}
 \wscal(h|_{T\Sigma},\phi_\Sigma+\eta_0)
 \ge\wscal(h,\phi)-H_1(G)-2\ln(v/L_1)
                         +c\Lambda^2+\frac{\mathcal L}{4}-\delta.
\end{equation}
We have obtained a large scalar allowance near the marked set.  The next
construction uses it to move the residual map out of the target band
there, while paying for the graph volume created in doing so.

\paragraph{Two graph prescriptions and their costs.}
We describe the graph calculations first, and then give a complete region
on which the prescriptions can be solved.  The first prescription will
produce a function $f$ within two of a small multiple of $\ln u_\alpha$,
with $|df|\le1$ in its graph metric.  The coordinate prescriptions then
use this controlled gradient to amplify $\bar J$ where $f$ is high;
each correction adds at most $e$ to its coordinate gradient in the final
graph metric, in addition to the earlier perturbation of $\bar J$.
Their scalar gains will pay the exact graph volume factors.

For a current metric and
weight $(k,\varphi)$, a length-one graph of a function $f$ has
\[
 L_f=(1+|df|_k^2)^{1/2},\quad U_f=\frac{\nabla_k f}{L_f},\quad
 k_f=k+df\otimes df,\quad \varphi_f=\varphi-\ln L_f,
 \quad D_f=-\operatorname{div}_{k,\varphi}U_f.
\]
Lemma~\ref{lem:graph-identity} says
\begin{equation}\label{eq:degree-graph-identity}
 \wscal(k_f,\varphi_f)-\wscal(k,\varphi)
    =D_f^2-2U_f(D_f)+|\mathrm{II}_f|^2+|d\ln L_f|_{k_f}^2.
\end{equation}
Also
\begin{equation}\label{eq:degree-graph-elementary}
 \begin{gathered}
 |U_f|_k\le1,\qquad U_f(f)=L_f-L_f^{-1},\qquad
 |df|_{k_f}\le1,\\
 dV_{k_f}=L_f\,dV_k,\qquad L_f-L_f^{-1}\ge\ln L_f.
 \end{gathered}
\end{equation}
The last inequality holds for every $L_f\ge1$ by differentiation.  These
relations will make the scalar cost and the volume change use the same
factor $L_f$.

For the first graph put
\[
      f_0=\sigma\ln u_\alpha,\qquad C_f=B_0+\Lambda,
      \qquad D_f=C_f\tan(f_0-f),
\]
with the argument between its two poles.  Here $B_0\ge2$ is fixed and
$\sigma>0$ will be small.  Pole avoidance gives $|f-f_0|<\pi/2<2$.
Differentiating the prescribed $D_f$ gives the exact terms
\begin{equation}\label{eq:degree-first-graph-differentiation}
 \begin{split}
 D_f^2-2U_f(D_f)
  ={}&D_f^2-2D_fU_f(\ln C_f)
      -2(C_f+D_f^2/C_f)U_f(f_0)\\
     &+2(C_f+D_f^2/C_f)(L_f-L_f^{-1}).
 \end{split}
\end{equation}
The final term is at least $2\ln L_f$ by
\eqref{eq:degree-graph-elementary}.  For the derivative of $C_f$, Young's
inequality and \eqref{eq:degree-moments} give
\[
 2|D_fU_f(\ln C_f)|
   \le0.8D_f^2+1.25|d\ln C_f|^2
   \le0.8D_f^2+\frac{1.25}{8}\mathcal L.
\]
The $f_0$ term has absolute value at most
\[
     2\sigma C(C_f+D_f^2/C_f)(\Lambda+1).
\]
Since $(\Lambda+1)/C_f\le1$, a sufficiently small $\sigma$ spends less
than the remaining $0.2D_f^2$ on its square term.  Moreover
$C_f(\Lambda+1)\le C_{B_0}(1+\Lambda^2)$, so a further small choice of
$\sigma$ spends less than half the $c\Lambda^2$ in
\eqref{eq:degree-weight-gain} and makes the remaining constant loss as
small as prescribed.  The available $\mathcal L/4$ covers
$1.25\mathcal L/8$.  Combining \eqref{eq:degree-weight-gain} and
\eqref{eq:degree-graph-identity}, and discarding nonnegative terms, thus
gives
\begin{equation}\label{eq:degree-first-graph-gain}
 \wscal(k_f,\varphi_f)\ge\wscal(h,\phi)-H_1(G)-2\ln(v/L_1)
                         +c_1\Lambda^2+2\ln L_f-\delta_1,
\end{equation}
where $c_1>0$ is fixed, while $\delta_1$ can be made arbitrarily small.  This
calculation is used only where the current base metric is the original
induced metric for which \eqref{eq:degree-weight-gain} was proved.

For the remaining graphs, let $\zeta(f)$ be a smooth function valued in
$[0,1]$, zero for $f\le K$ and one for $f\ge K+1$, with derivative bound
independent of the threshold $K$.  Write
$(k^{(1)},\varphi^{(1)})=(k_f,\varphi_f)$.  For row $j=2,\ldots,d$, its
graph quantities use the metric and weight after all preceding rows:
\[
 \begin{gathered}
 L'_j=(1+|df'_j|_{k^{(j-1)}}^2)^{1/2},\qquad
 U'_j=\frac{\nabla_{k^{(j-1)}}f'_j}{L'_j},\\
 D'_j=-\operatorname{div}_{k^{(j-1)},\varphi^{(j-1)}}U'_j,\\
 k^{(j)}=k^{(j-1)}+df'_j\otimes df'_j,\qquad
 \varphi^{(j)}=\varphi^{(j-1)}-\ln L'_j.
 \end{gathered}
\]
Prescribe successively
\begin{equation}\label{eq:degree-row-graph}
  \ell_j=A\zeta(f)\bar J_j,\qquad
  D'_j=B'\tan\left(\frac{\ell_j-e f'_j}{w}\right),\qquad
  J_j^+=\bar J_j+e f'_j.
\end{equation}
The constants $A,w,e,B'$ are positive; their order of choice is given
below.  Put $h^+=k^{(d)}$, $\phi^+=\varphi^{(d)}$, and
\[
                         \mathcal W=L_f\prod_{j=2}^dL'_j.
\]
The successive volume identities give $dV_{h^+}=\mathcal W\,dV_k$
relative to the current base metric $k$.  On the eventual child that
base will be the original induced metric, so this product will be the
density in \eqref{eq:degree-child-scalar}.  In the metric after the first graph, $|df|\le1$ by
\eqref{eq:degree-graph-elementary}, and later graph metrics only
increase.  The values and first derivatives of $\bar J$ are uniformly
bounded in the region in question.  It follows that in each current
metric
\begin{equation}\label{eq:degree-ell-bound}
             |d\ell_j|\le C_\ell\mathbf1_{\{f\ge K\}},
\end{equation}
where $C_\ell$ depends on $A$ and the earlier band bounds but not on
$K$, $M_*$, $k'$, or the final cover budget.

For a row graph the analogue of
\eqref{eq:degree-first-graph-differentiation} is
\[
 (D'_j)^2-2U'_j(D'_j)
  =(D'_j)^2-\frac2w\left(B'+\frac{(D'_j)^2}{B'}\right)U'_j(\ell_j)
       +\frac{2e}{w}\left(B'+\frac{(D'_j)^2}{B'}\right)
                      (L'_j-(L'_j)^{-1}).
\]
Choose $B'$ after $w,e$ so that
\begin{equation}\label{eq:degree-row-B-choice}
              \frac{B'e}{w}\ge2,\qquad
              \frac{2C_\ell}{wB'}\le\frac12.
\end{equation}
Then the differentiated $\ell_j$ spends at most half of $(D'_j)^2$ and
the fixed amount $2B'C_\ell/w$ on $\{f\ge K\}$, while the last term pays
$2\ln L'_j$ everywhere.  The other terms of
\eqref{eq:degree-graph-identity} are nonnegative.  After all $p$ rows,
the sum of fixed errors is therefore at most
$C_{\mathrm{rows}}\mathbf1_{\{f\ge K\}}$, with
$C_{\mathrm{rows}}$ independent of $K,M_*,k'$ and the cover budget.
On this support, $f<f_0+2$ implies
\[
  u_\alpha>\exp((K-2)/\sigma),\qquad
  \Lambda\ge c'\exp((K-2)/(\sigma\alpha))
\]
by \eqref{eq:degree-moments}.  Choose $K$ large enough that the remaining
$c_1\Lambda^2$ in \eqref{eq:degree-first-graph-gain} covers
$C_{\mathrm{rows}}$ there.  This proves the required payment for all row
graphs without any estimate on their individual volume expansions.

Pole avoidance in \eqref{eq:degree-row-graph} also gives the output
estimates that will be used for the band sides and its gradients:
\begin{equation}\label{eq:degree-output-estimates}
 \begin{split}
 |J_j^+-(1+A\zeta(f))\bar J_j|&<2w,\\
 |dJ_j^+|_{h^+}&\le |dG_j|_h+\varepsilon_{\mathrm{grad}},\\
 |J_j^+|&\ge |G_j|-\varepsilon_{\mathrm{val}}.
 \end{split}
\end{equation}
For the gradient inequality,
$|df'_j|\le1$ in its own graph metric and hence in $h^+$, and all these
metrics dominate the original induced metric and the parent metric on
tangent vectors.  Thus
$|dJ_j^+|_{h^+}\le|d\bar J_j|_h+e$.  For the value inequality use
$1+A\zeta\ge1$ in the first line.  The quantities
$\varepsilon_{\mathrm{grad}}$ and $\varepsilon_{\mathrm{val}}$ can be as
small as prescribed by the first-derivative and value perturbation
budgets, respectively, and by $e,w$.  We choose them to satisfy the
target budgets \eqref{eq:degree-tolerances} and also to be smaller than
all the fixed side margins used below.

The full dependency order is
\[
 \begin{aligned}
 &\text{perturbation budgets}\ \longrightarrow\ (\bar J,a_0,k)
       \ \longrightarrow\ R\\
 &\qquad\longrightarrow\ (\text{scalar error},\varepsilon_{\mathrm{base}})
       \ \longrightarrow\ B_0\ \longrightarrow\ \sigma\\
 &\qquad\longrightarrow\ A\ \longrightarrow\ (w,e)
       \ \longrightarrow\ B'\ \longrightarrow\ K\ \longrightarrow\ M_*\\
 &\qquad\longrightarrow\ (\text{cover budget},k',\text{radius lower bound}).
 \end{aligned}
\]
Each later choice may use the preceding ones.  The intended scalar error
is much smaller than $1/q$, and $\varepsilon_{\mathrm{base}}$ is small
enough for \eqref{eq:degree-potential-laplacian}.  The constants in
\eqref{eq:degree-moments} and the positive constant in
\eqref{eq:degree-weight-gain} are independent of the future cover, so
$B_0$ and $\sigma$ can be fixed for
\eqref{eq:degree-first-graph-gain} at this stage.  Choose $A$ large compared with
$T_p/a_0$, then $w$ small compared with $a_0$ and all value margins,
and $e$ small for \eqref{eq:degree-tolerances}.  Choose $B'$ by
\eqref{eq:degree-row-B-choice}, and only then $K$ to cover the row errors.
Also make $K$ larger than the baseline value
$\sigma\ln(\varepsilon_{\mathrm{base}}^\alpha)$ by a fixed margin.
Choose $M_*$ so large that
\eqref{eq:degree-potential-on-bad} implies
\begin{equation}\label{eq:degree-high-on-bad}
                        f_0>K+100\quad\hbox{on }D_{k'}.
\end{equation}
Finally choose the cover budget sufficiently small for all the active-mass
and error requirements, and obtain $k'\ge k$ and a cover with its positive
lower radius bound.  This final selection is possible for every fixed
$M_*$ by \eqref{eq:degree-cover-budget}.  It also provides all finite
uniform derivative bounds needed for the actual equations.  In
particular, none of the constants that forced the choices of $K$ or $M_*$
was allowed to depend on the final level or cover.

\paragraph{Solving on a complete regular region.}
The thresholds have distinct jobs: $K+100$ places the marked set beyond
the auxiliary cutoff,
$K+12$ bounds the auxiliary region for solving, and $K+7$ bounds the
child inside the part where its base metric is unchanged.  The lower
threshold $K+4$ will already put the output outside the target near the
potential cutoff side.  Their gaps provide the spatial buffers below
through the uniform derivative bounds for $f_0$.

$\Sigma_s$ need not be complete.  Use the operations
only on the open region
\begin{equation}\label{eq:degree-auxiliary-region}
       \Omega=\Sigma_s\cap\{|G|_\infty<t_2,\ f_0<K+12\}.
\end{equation}
By \eqref{eq:degree-high-on-bad} it avoids $D_{k'}$ with a large margin in
$f_0$; by $t_2<t_3$ it also has a margin within the region of the tests.
The uniform derivative bounds on $f_0$ turn these into uniform spatial
margins.  Thus the graph tests through level $k'$ give regular graph
charts with uniform buffers and all estimates on a neighborhood of this
region.  The estimates for derivatives of $\ln v$ apply there.  The
functions $u_j,\Lambda,f_0,\bar J$ have uniform derivatives because of
the radius lower bound.  The weight
$\phi_\Sigma+(\ln u_\alpha+\ln u_{\alpha+4})/4$ therefore has the
positive-order derivative bounds required for a graph solve.

Choose a smooth function $0<\xi\le1$ on $\Omega$ which tends to zero
at the sides in \eqref{eq:degree-auxiliary-region} and is one, with
buffers, through the smaller bounds $|G|_\infty\le t_1$ and
$f_0\le K+7$.  It can be made from smooth functions of
$t_2^2-G_i^2$ and $K+12-f_0$, linear near zero and cut off to one away
from zero.  Its derivatives in the original regular charts are bounded.
Replace the induced metric on $\Omega$ by
\[
                              k_{\mathrm{aux}}=\xi^{-2}h|_{T\Sigma}.
\]
This metric is complete with uniform buffered bounded geometry.  Here is
the local verification, which also accounts for possible ends of the
original leaves.  A ball in an original regular chart of radius a small
fixed multiple of $\xi(x)$ remains in $\Omega$, since $\xi(x)$ is bounded
by a constant times each scalar distance to a side.  The bound on
$d\xi$ makes $\xi$ comparable to $\xi(x)$ on that ball.  Rescale the
coordinate ball by $\xi(x)$.  The coefficients of $\xi^{-2}h$ and their
inverse then have uniform bounds of every order: for the $m$-th derivative
of the conformal factor the new terms are products of ratios bounded by
constants times $\xi(x)^m|\nabla^m\xi|/\xi$, and these are bounded for
$m\ge1$.
The rescaled boxes have uniform inner buffers.  They also prove
completeness.  A curve of finite remaining length, or a Cauchy sequence,
eventually stays within the outer part of one such buffered box and
converges there; it cannot escape a leaf through a side or an uncharted
end.  Equivalently, approach to a side already has infinite length
because $|d\ln\xi|_{k_{\mathrm{aux}}}$ is bounded.  This argument used
only the controlled charts near \eqref{eq:degree-auxiliary-region},
not completeness of a parent leaf.

The restricted chart descriptions are Borel and accessible charts still
have a countable Borel enumeration.  The auxiliary atlas may use further
countable subdivisions; its total transverse measure is irrelevant for
the graph existence theorem.  Apply Lemma~\ref{lem:graph-prescription}
on these complete leaves.  For the first prescription its positive
coefficient is $B_0+\Lambda$, its height target is $f_0$, and its
constants $e,w$ in that lemma are both one.  For a row prescription the
coefficient is the positive constant $B'$, the target is $\ell_j$, and
the constants are the chosen $e,w>0$.  All coefficients have the
uniform smooth bounds just established.  That lemma gives the bounded
solutions with a uniform margin from the poles and all derivative
bounds, and gives their Borel dependence from uniqueness in the
countably enumerated charts.  It applies successively: a solved graph
metric dominates the complete base, is complete, and has uniform
geometry by the solution bounds; its changed weight again has uniform
positive-order derivative bounds.  In applying the scalar estimates
\eqref{eq:degree-weight-gain}--\eqref{eq:degree-first-graph-gain}, use
only the smaller region where $\xi=1$.  There the auxiliary metric and
the original induced metric agree with all derivatives.

Define the child to be the open measured band
\begin{equation}\label{eq:degree-child-domain}
 X_s^+=\{x\in\Omega: |G|_\infty<t_1,\ f_0<K+7,
                                      \ |J^+|_\infty<T_p\},
\end{equation}
with the final graph metric $h^+$, final graph weight $\phi^+$, and
underlying boundary orientation.  Its artificial sides stay away from
each inner residual subcube.  For the $|G|=t_1$ sides, confinement gives
$|G_1|<\tau+b<T_p$, so any approach to such a side has a residual row
$|G_j|$ near $t_1$.  Choose
$\varepsilon_{\mathrm{val}}<(t_1-T_p)/4$; then
\eqref{eq:degree-output-estimates} puts $|J_j^+|$ beyond $T_p$ already
on a uniform neighborhood of that side.  For the $f_0=K+7$ side, if
$f_0\ge K+4$ some kernel is active, since the baseline value of $f_0$
is less than $K$.  Its support lies in $|\bar J|_\infty>3a_0$.
Moreover $f>f_0-2\ge K+2$, so $\zeta(f)=1$.  Choose $A$ so that
$(1+A)3a_0>T_p+4a_0$ and take $w<a_0/2$.  The first line of
\eqref{eq:degree-output-estimates} then puts some $|J_j^+|$ beyond
$T_p$ with a positive uniform margin.  This occurs before reaching
$f_0=K+7$.  In particular the entire child lies in the region where
$\xi=1$, with buffers for all its inner subcubes.

These observations verify every part of
Definition~\ref{def:degree-band} for the child.  To be explicit about its
atlas, use the recognized regular graph patches through level $k'$ in
each parent plaque, restricting their transversals as necessary.  There
are only boundedly many such tests per plaque.  The derivative bounds
for $G,f_0,J^+$ and the side margins just proved permit a further bounded
subdivision covering any chosen inner residual subcube with buffers.
The solved graph metrics and weights have uniform derivatives in these
charts.  Holonomy is inherited from the parent overlap maps, with the
finite patch indices when needed.  Hence both the total transverse
measure and the plaque overlap bound persist.  The transverse measure
is precisely the restricted parent measure.  It has no lapse factor and
no graph-volume factor; those belong to the leafwise integrations.

\paragraph{Preserving the degree through the modifications.}
Use in \eqref{eq:degree-boundary-degree} a residual form $\beta$ supported
in $|z|_\infty<a_0$, reducing $a_0$ if necessary.  First homotope $J$ to
$\bar J$ on the whole parent boundary current.  The small value
perturbation ensures that the preimages of this support throughout the
homotopy remain in a fixed old inner band.  One can justify invariance
directly by the homotopy formula.  The difference of the two pullbacks
of $\beta$ is $d\mathcal H\beta$, where $\mathcal H\beta$ is its usual
transgression along the homotopy.  Multiply by the cutoff
$\chi_{\mathrm{deg}}(G_1)$ used
in \eqref{eq:degree-boundary-degree}.  The current $\partial[E_s]$ is
closed, and the term differentiating this cutoff meets no boundary on
the small residual support, by confinement.  Measured current Stokes,
with the local mass bounds already proved, therefore makes the integral
of the difference zero.  This uses the boundary current across its
singular set and so does not assume regularity of the homotopy there.

The resulting integral for $\bar J$ is entirely in the regular region
used for the solves.  Indeed $|\bar J|_\infty<a_0$ avoids a neighborhood
of $D_k$ by \eqref{eq:degree-gap} and avoids every kernel support, which
lies in $|\bar J|_\infty>3a_0$.  Thus $f_0$ has its baseline value there.
The small perturbation also keeps $|G|$ in the smaller inner band.
Now homotope on this regular region first from $\bar J$ to
$(1+A\zeta(f))\bar J$ and then to $J^+$.  Along the first homotopy the
multiplier is at least one, so a preimage of the small target support
has $|\bar J|_\infty<a_0$.  Along the second, the first line of
\eqref{eq:degree-output-estimates} shows that it has
$|\bar J|_\infty<a_0+2w<2a_0$.  Both supports thus avoid the marked
region and the kernel supports, and retain the same spatial and
$G_1$ margins.  Ordinary measured Stokes on the regular buffered charts
applies to these homotopies.  The final support lies in
\eqref{eq:degree-child-domain}, so its integral is the measured degree
of the child.  It equals $\mathcal D_X(G)$.

On the child the initial auxiliary metric is the induced metric, so
\eqref{eq:degree-graph-elementary} verifies the volume identity in
\eqref{eq:degree-child-scalar} for the product $\mathcal W$ defined above.
Choose all the adjustable scalar errors above with sum at most $1/q$.
The absorption of the row errors and
\eqref{eq:degree-first-graph-gain} give the scalar inequality in
\eqref{eq:degree-child-scalar} on the child.
This completes the construction in high dimensions.  In the case
$d\le7$ treated above, use $\mathcal W=1$.  The volume identity is then
immediate, and \eqref{eq:degree-slice-bound} gives the scalar inequality
in \eqref{eq:degree-child-scalar}.  In dimension $p=0$ the same statements
hold on the counted points with the zero-dimensional scalar convention.

\begin{proof}[Completion of the proof of Lemma~\ref{lem:degree-descent}]
The coarea assertion is \eqref{eq:degree-coarea-proof}.  The constructions
above give a controlled child of degree $\mathcal D_X(G)$ for each typical
$s$.  Its coordinate bound follows from
\eqref{eq:degree-output-estimates} and \eqref{eq:degree-tolerances}:
\[
 L_j|dJ_j^+|_{h^+}\le B-\frac dq+L_j\varepsilon_{\mathrm{grad}}
                         <B-\frac pq.
\]
For a low-dimensional child the same statement holds with zero error.
The horizon condition \eqref{eq:degree-horizon} for its rows is unchanged,
since $q,B,\tau,L_2,\ldots,L_d$ are unchanged.

Every point with $|J^+|_\infty<b_p$ comes from
$|G|_\infty<b_d$.  For residual rows this follows from the last line of
\eqref{eq:degree-output-estimates} and
$\varepsilon_{\mathrm{val}}<b_d-b_p$; in low dimensions it is immediate.
For the first row confinement gives
\[
       |G_1|<\tau+b<\tau+\frac{\tau}{10q}<b_d.
\]
The same observation covers the counted case $p=0$.

The function $x\mapsto(x-1)_+$ is one-Lipschitz, so
\eqref{eq:degree-output-estimates} gives
\[
            \sum_{j=2}^dH_j^+\le\sum_{j=2}^dH_j(G)
                         +2\sum_{j=2}^dL_j\varepsilon_{\mathrm{grad}}.
\]
Use this and the scalar and volume relations in
\eqref{eq:degree-child-scalar} to compare the integrands on the child:
\begin{equation}\label{eq:degree-integrand-cancellation}
 \begin{split}
 &\exp\left(\frac{-\wscal(h^+,\phi^+)+\sum_{j=2}^dH_j^+}{2}\right)
                  \mathcal W\\
 &\quad\le
 \exp\left(\frac1{2q}+\sum_{j=2}^dL_j\varepsilon_{\mathrm{grad}}\right)
       \frac v{L_1}
       \exp\left(\frac{-\wscal(h,\phi)+\sum_{i=1}^dH_i(G)}2\right)\\
 &\quad\le e^{3/(5q)}\frac v{L_1}\Theta.
 \end{split}
\end{equation}
The factor $\mathcal W$ cancels exactly against the term
$2\ln\mathcal W$ in the scalar bound.  The factor $v/L_1$ is the lapse
factor in coarea.  This calculation makes no bound on either factor.
For a low-dimensional child use \eqref{eq:degree-slice-bound}, with no
row error; for $p=0$ the same estimate is the pointwise bound for a
counted point.  Integrate \eqref{eq:degree-integrand-cancellation} over
the child and use its containment in the parent $b_d$ region.  This is
\eqref{eq:degree-child-integral}, for example with $C_*=1$.
\end{proof}

\subsection{Completion of the induction}

\begin{proof}[Proof of Theorem~\ref{thm:degree}]
The case $d=0$ is the signed-count inequality in the statement.  Suppose
the theorem is known in dimension $p=d-1$, with the same $q,B,\tau$ and
the residual lengths.  Its hypotheses hold for the child of
Lemma~\ref{lem:degree-descent}, including the unchanged strict horizon
condition.  For almost every $s$ the induction hypothesis and
\eqref{eq:degree-child-integral} give
\[
 |\mathcal D_X(G)|(2\tau)^p
 \le \frac{e^{C_1p/q+C_*/q}}{\prod_{i=1}^dL_i}
          \int_{\Sigma_s\cap\{|G|_\infty<b_d\}}v_s\Theta\,dA_h.
\]
For $p=0$ this is the same use of the unsigned-count bound.  Take
$C_1\ge C_*$.  Integrate this scalar inequality over $-\tau<s<\tau$
and apply \eqref{eq:degree-coarea} with
$F=\mathbf1_{\{|G|_\infty<b_d\}}\Theta$.  The factor $2\tau$ on the
left changes its power from $p$ to $d$, and the right is bounded by
the integral in \eqref{eq:degree-estimate}.  This proves the induction.

There is no implicit measurable selection of children in this last step.
For each typical $s$ a child proves the displayed scalar inequality
against the same parent integral.  That parent integral is a measurable
function of $s$ by the lapse construction and is governed by
\eqref{eq:degree-coarea}; it is the only function being integrated.
Its integral is finite because $\Theta$ is bounded on the controlled
$b_d$ region, which has finite measured volume.
The constants $C_0$ and $C_1$ are absolute: $C_0$ came from the fixed
pressure profile, and each descent costs at most $e^{C_*/q}$ after the
arbitrarily small slice errors were chosen.  All geometry-dependent
constants were used only to make those choices.
\end{proof}
